% Lösungen Einheit 3 – Grundwert berechnen
\einheitenkopf*{Einheit 3 von 4 · Grundwert berechnen}
\erg{31}{a) Prozentwert \quad b) Prozentsatz \quad c) Grundwert \quad d) Prozentsatz \quad e) Grundwert}
\erg{32}{Beispiel 60\,kg \quad a) 1 Kästchen $=$ 5\,€, 50\,€ \quad b) 1 Kästchen $=$ 7\,m, 70\,m \quad c) 1 Kästchen $=$ 2,60\,kg, 26\,kg \quad d) 1 Kästchen $=$ 11\,€, 44\,€}
\erg{33}{Beispiel 42\,€ \quad a) $9 \cdot 2 = 18$\,€ \quad b) $9 \cdot 4 = 36$\,€ \quad c) $9 \cdot 5 = 45$\,€ \quad d) $9 \cdot 10 = 90$\,€ \quad e) Der Teil (9\,€) bleibt gleich; je kleiner die Prozentzahl, desto größer ist das Ganze.}
\erg{34}{Beispiel 400\,kg \quad a) 1\,\% $=$ 3\,€, 100\,\% $=$ 300\,€ \quad b) 1\,\% $= 2{,}6$\,kg, 100\,\% $=$ 260\,kg \quad c) 75\,\% $=$ 18 Kinder, 25\,\% $=$ 6 Kinder, 100\,\% $=$ 24 Kinder}
\erg{35}{a) Prozentwert: $0{,}3 \cdot 80 = 24$\,kg \quad b) Prozentsatz: $75 : 250 = 0{,}3 = 30\,\%$ \quad c) Grundwert: 1\,\% $= 24 : 40 = 0{,}6$, 100\,\% $=$ 60 Kekse \quad d) 1\,\% $= 12 : 15 = 0{,}80$\,€, 100\,\% $=$ 80\,€ (nicht 15\,\% von 12\,€)}
\erg{36}{a) Sara rechnet 25\,\% vom Teil; der Teil ist aber schon 25\,\%, das Ganze muss größer sein: $12 \cdot 4 = 48$\,€ \quad b) 1\,\% $=$ 0,35\,kg, 100\,\% $=$ 35\,kg}
\erg{37}{a) Das Ganze sind 100\,\%, also 100-mal so viel wie 1\,\%. \quad b) Stimmt: 150\,\% sind mehr als das Ganze. 1\,\% $=$ 0,20\,€, 100\,\% $=$ 20\,€.}
\erg{38}{a) 1\,\% $= 42 : 70 = 0{,}60$\,€, 100\,\% $=$ 60\,€ \quad b) $60 - 42 = 18$\,€ \quad c) Nein: In 4 Wochen spart sie 16\,€, es fehlen 18\,€. Nach 5 Wochen reicht es.}
